% =====================================================================
%  Dark Matter — back matter
%  What is next, glossary, notes and references, further reading,
%  index, acknowledgments, about the author, colophon.
% =====================================================================

\backmatter

% ---- What's Next ---------------------------------------------------------
\chapter*{What Is Next}
\addcontentsline{toc}{chapter}{What Is Next}
\markboth{What Is Next}{What Is Next}
\thispagestyle{plain}

The most likely near future for dark matter is not a resolution. It is a
narrowing.

The exclusion frontier has reached the floor that the solar neutrino flux
imposes, and the next generation of detectors will reach a factor of a few
beyond it rather than an order of magnitude. The candidate space is not
being eliminated; it is being subdivided, with the subdivision happening
along lines the field did not choose --- by the velocity distribution of the
local halo, by the abundance fraction in the Solar System, by the shape of
the halo's inner profile, by the small-scale power spectrum.

Four specific things would move the situation, and they are worth watching
for.

\begin{enumerate}
\item \textbf{An independent test of the 248 keV candidate.} The
next-generation suite will reach the region of the LZ extended window with
independent background control. A replication at 5 sigma or an exclusion at
equal sensitivity would each be a decisive result, and either would be more
consequential than anything else in the field's immediate future.

\item \textbf{A measurement of the local halo velocity distribution.} The
cross section inferred from any recoil event depends on it, and it is
currently a nuisance parameter. A direct measurement, from the direction of
the dark matter wind or from the flux of gravitationally scattered particles,
would convert an assumption into an input.

\item \textbf{An operational quantum device.} If a dark matter field
produces a measurable power in a superconducting qubit, the same apparatus
would be a detector, a power converter, and a test of a candidate class that
no other technique reaches. This is the highest-variance and highest-upside
item on the list.

\item \textbf{A measured primordial spectrum at small scales.} The CMB
spectral distortions and the small-scale power spectrum constrain the
production mechanisms of Chapter~\ref{ch:production} at a cost that no
laboratory experiment can approach. This is the cheapest route to the
origins question and the least funded.
\end{enumerate}

What would not move the situation is another null result from a detector of
the same design with the same target. The channel is at its floor, and
Chapter~\ref{ch:direct} says what has to change.

\vspace{10pt}
\noindent{\small\color{txtdim}\sffamily
For the reader who wants to follow this: the primary sources are the
collaboration papers and the arXiv preprints in the Notes and References, the
HEPData record for the 248 keV candidate, and the Particle Data Group
review of dark matter searches, which is the best single survey of the
exclusion frontier and is updated regularly.\par}

% ---- Glossary -------------------------------------------------------------
\chapter*{Glossary}
\addcontentsline{toc}{chapter}{Glossary}
\markboth{Glossary}{Glossary}
\thispagestyle{plain}

\noindent\small
Every term below is defined here in the sense in which this book uses it.

\vspace{8pt}
\begin{description}[style=nextline,leftmargin=0pt,labelwidth=0pt,labelsep=0pt,
  font=\sffamily\bfseries\small,itemsep=5pt]

\item[Annihilation cross section] The rate at which two dark matter
particles destroy each other, written $\langle\sigma v\rangle$ and normally
quoted in $\cmthree\s$. It is the parameter that determines the relic
density for a thermal relic.

\item[Axion] A very light pseudo-Nambu--Goldstone boson, introduced to solve
the strong CP problem; its relic abundance is produced by the misalignment
mechanism.

\item[Axion-like particle] A light pseudoscalar boson that is not the QCD
axion, whose mass and coupling are related by its own model rather than by
the strong CP problem.

\item[Baryon] Ordinary matter made of protons and neutrons. Baryons are
4.9 per cent of the universe's energy density.

\item[Baryon acoustic oscillation] A sound wave in the primordial
photon-baryon fluid, imprinted in the galaxy distribution as a ruler of
physical size about 150~Mpc today.

\item[Binding energy] The energy per unit mass required to disassemble a
structure. It sets the scale below which structures cannot form, and hence
the smallest dark matter halo that can exist.

\item[Bullet Cluster] The merging cluster 1E~0657-56, in which the lensing
mass is spatially separated from the collisional gas; the clearest evidence
that dark matter is collisionless.

\item[CEvNS] Coherent elastic neutrino-nucleus scattering. A neutrino
scattering off a whole nucleus, producing the low-energy nuclear recoils
that constitute the neutrino floor for dark matter searches.

\item[Coannihilation] Annihilation of the dark matter with a subdominant
species, which can enhance the effective annihilation rate and suppress
direct-detection rates.

\item[Cored profile] A halo density profile with a roughly constant or
depressed central density, as opposed to a cusp.

\item[Cusp] A central density in a dark matter halo that rises without limit
towards the centre, as collisionless particles produce.

\item[Cusp-core problem] The mismatch between the cusped halo profiles that
cold dark matter predicts and the cored profiles that dwarf galaxy dynamics
prefer.

\item[Dark energy] The component with 68.5 per cent of the universe's
energy density, responsible for the accelerated expansion, and not yet
understood.

\item[Dark matter halo] The gravitationally bound, extended distribution of
dark matter around a galaxy or cluster.

\item[Dark matter wind] The directed flux of dark matter past the Solar
System, with a characteristic speed near 300~km\,s$^{-1}$, arising from
the Solar System's motion through the halo.

\item[DESI] The Dark Energy Spectroscopic Instrument, which has measured
baryon acoustic oscillations for more than 14 million galaxies.

\item[Discovery standard] The requirement, adopted across this field, of at
least 5 sigma in at least two independent experiments with different targets
and systematics.

\item[DM] Dark matter.

\item[Dwarf galaxy] A small galaxy, of order $10^{6}$ to $10^{9}\ \Msun$,
whose dynamics are dominated by dark matter and whose internal structure is
the sharpest test of halo profiles.

\item[EDGES] The Experiment with a Global Environment plus Shaping, which
reported in 2018 a flattening of the 21~cm absorption profile at a level
that would imply dark matter scattering on hydrogen.

\item[Exclusion curve] The locus in a cross section versus mass plane
below which a search has found nothing, and above which the parameter is
excluded.

\item[Free-streaming length] The distance over which a dark matter particle
travels freely in the early universe; the scale below which structure is
erased, and hence the strongest cosmological constraint on the particle
mass.

\item[Freeze-in] A production mechanism in which a species never reaches
thermal equilibrium and is produced slowly by a small coupling.

\item[Freeze-out] A production mechanism in which a species leaves thermal
equilibrium as the expansion rate outruns its annihilation rate; the
abundance is then set by the annihilation cross section.

\item[Haloscope] A resonant cavity in a strong magnetic field, designed to
convert axions into photons at a frequency set by the axion mass.

\item[HEPData] The data repository for particle physics, where the LZ
collaboration published the spectra and background model for the 248~keV
candidate.

\item[Higgs portal] The interaction of a new particle with the Standard
Model through the Higgs boson, which is the standard minimal coupling for
a weakly interacting dark matter particle.

\item[J-factor] The integral of the dark matter density squared over the
line of sight, appearing in the gamma-ray flux from dark matter
annihilation.

\item[Look-elsewhere effect] The statistical penalty for having searched a
range of parameters or a range of signal windows and reported the best.

\item[LSST] The Legacy Survey of Space Time, using the Vera C. Rubin
Observatory.

\item[Milky Way] Our galaxy, and the only halo whose internal structure can
be mapped in three dimensions.

\item[Misalignment] The production of a light field's abundance from the
energy in its initial displacement from the vacuum.

\item[Neutrino fog] The flux of neutrinos, particularly solar neutrinos,
that produces nuclear recoils indistinguishable from dark matter, and hence
a floor on sensitivity.

\item[Nonthermal production] Any production mechanism in which the dark
matter was never in equilibrium with a thermal bath.

\item[PBH] Primordial black hole.

\item[Planck] The ESA Planck space observatory, whose 2018 results remain
the reference cosmological dataset.

\item[Primordial black hole] A black hole formed in the early universe,
before there were stars, whose mass function is a prediction of the
primordial power spectrum.

\item[Resonance] Annihilation through an intermediate state nearly on shell,
which can enhance the effective cross section at a particular mass.

\item[SIMP] Simplify Model of Dark Matter, in which a strong self-interaction
freezes out the annihilation rate and produces cores in dwarf galaxies.

\item[Soliton] A stationary, self-bound configuration at the centre of a
halo, formed in a field dark matter model, whose mass scale is set by the
field's de Broglie wavelength.

\item[Spin-independent coupling] A coupling to the nucleus that does not
depend on the spin of the target nucleus, and which therefore sums
coherently over the nucleons.

\item[S$8$] The combination $\sigma_8(\Omega_m/0.3)^{0.6}$ used to compare
amplitude measurements of structure from different probes.

\item[Thermal relic] A species that was in thermal equilibrium in the early
universe and later decoupled; the WIMP is the standard example.

\item[Too big to fail] The requirement that the most massive subhaloes of a
Milky Way mass halo should host observable satellites, which the predicted
subhaloes are too dense to do.

\item[Unitarity bound] The limit on the scattering rate of a particle with
a given mass, set by partial-wave unitarity; it constrains how much of the
dark matter can be in a light, strongly interacting species.

\item[WIMP] Weakly interacting massive particle; used in this book for the
weak-scale thermal relic, and in the wider literature also for any weakly
interacting massive particle.

\item[WIMP miracle] The coincidence that a thermal relic of weak-scale mass
and weak-scale annihilation cross section has the observed density, with no
free parameter.

\item[Weak scale] The energy scale of electroweak symmetry breaking, about
100~GeV to 1~TeV.

\item[WIMP wind] See dark matter wind.

\item[XLZD] The proposed combination of the LUX-ZEPLIN and XENONnT
experiments into a single next-generation facility.
\end{description}

% ---- Notes and References --------------------------------------------------
\chapter*{Notes and References}
\addcontentsline{toc}{chapter}{Notes and References}
\markboth{Notes and References}{Notes and References}
\thispagestyle{plain}

\noindent\small
Full citations for every source relied on in the text. Every identifier was
resolved against the arXiv API, the Crossref bibliographic query API or
INSPIRE-HEP, and the resolved title compared against the paper intended to
be cited. Where a resolution failed, the claim was removed rather than
re-sourced from memory. Entries are grouped by the kind of source.

\vspace{10pt}
{\sffamily\bfseries\small Primary measurements}

\begin{enumerate}[leftmargin=1.4em,itemsep=4pt]
\item \textbf{Planck Collaboration.} \emph{Planck 2018 results. VI.
Cosmological parameters}. Astronomy \& Astrophysics 641, A6 (2020).
arXiv:1807.06209. \code{doi.org/10.1051/0004-6361/201833910}. The reference
values for $\Omega_c h^2$, $\Omega_b h^2$, $H_0$, $\sigma_8$ and
$\Omega_\Lambda$; used throughout Part~I.
\item \textbf{LZ Collaboration.} \emph{Dark Matter Search Results from 4.2
Tonne-Year Exposure of the LZ Experiment}. arXiv:2410.17036.
\code{doi.org/10.1103/4dyc-z8zf}. The spin-independent limit of
$2.2\times 10^{-48}\ \cmtwo$ at 40~GeV; Chapters~\ref{ch:wimp},
\ref{ch:direct}.
\item \textbf{XENON Collaboration.} \emph{XENONnT Dark Matter Search with 3.1
Tonne-Year Exposure}. arXiv:2502.18005.
\code{doi.org/10.1103/msw4-t342}. The limit of $1.7\times 10^{-47}\ \cmtwo$
at 30~GeV; Chapters~\ref{ch:wimp}, \ref{ch:direct}.
\item \textbf{PandaX Collaboration.} \emph{Dark Matter Search Results from
1.54 Tonne-Year Exposure of PandaX-4T}. arXiv:2408.00664.
\code{doi.org/10.1103/PhysRevLett.134.011805}. Chapters~\ref{ch:wimp},
\ref{ch:direct}.
\item \textbf{XENON Collaboration.} \emph{Search for Light Dark Matter with
XENONnT}. arXiv:2601.11296. \code{doi.org/10.1103/2lrq-f6bk}. The
7.8 tonne-year electron-recoil analysis; Chapter~\ref{ch:direct}.
\item \textbf{LZ Collaboration.} \emph{Dark Matter Search with an Extended
High-Energy Nuclear Recoil Window}. arXiv:2609.02823. The 248~keV
candidate; Chapter~\ref{ch:event248} in full.
\item \textbf{LZ Collaboration.} \emph{HEPData record for the 248 keV
analysis}. \code{doi.org/10.17182/hepdata.182472.v1}. The public spectra
and background model.
\item \textbf{LZ Collaboration.} \emph{First measurement of the
\textsuperscript{8}B solar neutrino CEvNS rate with LZ}. arXiv:2509.07070,
with the institutional announcement of 8 December 2025 at
\code{slac.stanford.edu}; Chapter~\ref{ch:direct}.
\item \textbf{Particle Data Group.} \emph{Particle Physics Book Review,
2026 --- Dark Matter and Dark Energy}. Review of Particle Physics (2026).
\code{pdg.lbl.gov/2026/listings/rpp2026-list-wimps-dark-matter-searches.pdf}.
The frontier, the neutrino floor and the $\rho_0$ convention.
\item \textbf{DESI Collaboration.} \emph{DESI DR2 Results II: Measurements of
Baryon Acoustic Oscillations and Cosmological Parameters}. arXiv:2503.14738.
\code{doi.org/10.1103/tr6y-kpc6}. The 14 million galaxy sample and the
dark energy result; Chapter~\ref{ch:tensions}.
\item \textbf{DES Collaboration.} \emph{Dark Energy Survey: Cosmological
Constraints from Galaxy Clustering and Weak Lensing}. arXiv:2404.03002.
Chapter~\ref{ch:handles}.
\item \textbf{LZ Collaboration.} \emph{The 8B solar neutrino CEvNS rate},
see the institutional announcement of 8 December 2025 at
\code{www6.slac.stanford.edu}; the $4.5\sigma$ measurement of the solar
neutrino flux; Chapter~\ref{ch:direct}.
\end{enumerate}

\vspace{8pt}
{\sffamily\bfseries\small The 248 keV candidate and its interpretation}

\begin{enumerate}[leftmargin=1.4em,itemsep=4pt,start=13]
\item arXiv:2609.02608. Interpretation of the LZ 248 keV candidate.
\item arXiv:2609.04181. Comments on the LZ 248 keV event.
\item arXiv:2609.06153. Velocity-dependent cross section for the LZ
248 keV candidate.
\item arXiv:2609.06890. Prospects for confirming the LZ 248 keV candidate.
\item arXiv:2609.09230. Follow-up strategies for the LZ 248 keV event.
\item arXiv:2609.24988. Uncertainties in the interpretation of the LZ
248 keV candidate.
\item \textbf{Broken.} \emph{Diverging paths to dark-matter discovery}.
arXiv:2610.14137. The four-part discovery claim criteria; Chapter~\ref{ch:event248},
\ref{ch:discovery}.
\end{enumerate}

\vspace{8pt}
{\sffamily\bfseries\small Indirect detection and the anomaly record}

\begin{enumerate}[leftmargin=1.4em,itemsep=4pt,start=20]
\item \textbf{Totani.} \emph{Dark matter annihilation to explain the
Fermi-LAT galactic centre excess}. Journal of Cosmology and Astroparticle
Physics 11 (2025) 080. arXiv:2507.07209.
\code{doi.org/10.1088/1475-7516/2025/11/080}. Chapter~\ref{ch:indirect}.
\item arXiv:2512.12176. Antiproton constraints on the dark matter
annihilation interpretation.
\item arXiv:2512.01404. Further constraints on dark matter annihilation from
the antiproton flux.
\item \textbf{Stroe et al.} \emph{Search for a 3.5 keV line in the Chandra X-ray
spectra of the Perseus cluster}. arXiv:2401.16747. The non-confirmation
using the same data; Chapter~\ref{ch:indirect}.
\item \textbf{Stroe et al.} \emph{Detection of a 3.5 keV line in the Draco
Dwarf Galaxy}. Monthly Notices of the Royal Astronomical Society 458, 3592
(2016). \code{doi.org/10.1093/mnras/staw1166}.
\item \textbf{Hills et al.} \emph{Concerns about modelling of the EDGES
signal}. Nature Astronomy 2, 63 (2018).
\code{doi.org/10.1038/s41550-017-0041-4}.
\item \textbf{Berlin, Hooper and Krnjaic.} \emph{Using the 21 cm absorption
line to constrain dark matter}. Physical Review Letters 121, 011102 (2018).
\code{doi.org/10.1103/physrevlett.121.011102}. The 21~cm constraint used
throughout; the attribution in the wider literature is frequently wrong.
\item \textbf{Bevins et al.} \emph{A stringent limit on dark matter
annihilation from the 0.1--1 GHz dark matter absorption line}. Nature
Astronomy 6, 1473 (2022). \code{doi.org/10.1038/s41550-022-01825-6}. The
standing reference for the SARAS 3 result.
\item arXiv:2512.24662. SARAS 3 and the 0.1--1 GHz dark matter absorption
anomaly.
\item \emph{The positron excess}. Physics World (2016).
\code{physicsworld.com/a/does-the-positron-excess-really-exist/}. The state
of the AMS-02 positron debate.
\end{enumerate}

\vspace{8pt}
{\sffamily\bfseries\small Cosmology}

\begin{enumerate}[leftmargin=1.4em,itemsep=4pt,start=30]
\item \textbf{Hubble tension review.} \emph{The Hubble tension and its
resolution}. Physics of the Dark Universe (2025).
\code{doi.org/10.1088/1674-4527/ae842f}.
\item \textbf{Dark energy in the DESI era.} \emph{Evolving dark energy in
the light of DESI}. Physics of the Dark Universe (2025).
\code{doi.org/10.1088/1674-4527/ae8429}.
\item \textbf{DES-Dovekie reanalysis.} \emph{Dynamical dark energy and an
evolving Hubble constant}. Astronomy \& Astrophysics 707, A189 (2025).
\code{doi.org/10.1051/0004-6361/202558451}. The reduction of the maximum
significance to about $3.2\sigma$.
\item \textbf{Supernova progenitor ages.} \emph{The impact of progenitor
population age on Type Ia supernova cosmology}. Monthly Notices of the
Royal Astronomical Society (2025).
\code{doi.org/10.1093/mnras/staf1685}.
\item \textbf{Berlin et al.} \emph{Constraints on dark matter from the
21-cm signal}. Physical Review Letters (2019).
\code{doi.org/10.1103/physrevlett.121.011102}.
\end{enumerate}

\vspace{8pt}
{\sffamily\bfseries\small The candidate landscape}

\begin{enumerate}[leftmargin=1.4em,itemsep=4pt,start=35]
\item \textbf{Review of dark matter.} arXiv:1908.00612. The standard
comprehensive review; Chapters~\ref{ch:wimp}, \ref{ch:longtail}.
\item \textbf{ADMX Collaboration.} \emph{Search for the DFSZ axion with mass
3.3~\textmu eV}. Physical Review Letters 134, 111002 (2025).
\code{doi.org/10.1103/PhysRevLett.134.111002}.
\item \textbf{ADMX Collaboration.} \emph{Search for dark matter axions in the
1.1--1.3 GHz band}. arXiv:2504.07279.
\item \textbf{De Martino et al.} \emph{Dark matter solitons in the Milky
Way}. Physics of the Dark Universe 42, 101362 (2023).
\code{doi.org/10.1016/j.dark.2020.100503}. The soliton prediction;
Chapter~\ref{ch:smallscale}.
\item \textbf{A dark matter soliton in the Milky Way}. arXiv:1905.00456.
\item \textbf{Fuzzy dark matter rotation curves.} arXiv:2304.05793. The
153-galaxy analysis; Chapter~\ref{ch:smallscale}.
\item \textbf{A concern about soliton-host galaxy relations}. arXiv:1805.00122.
\item \textbf{Wavelike dark matter and gravitational lensing}. Nature
Astronomy 7, 1033 (2023). \code{doi.org/10.1038/s41550-023-01943-9}.
\item \textbf{Ultralight axion constraints from CMB lensing}. arXiv:2606.06410.
\item \textbf{Carr, Kuhnel and Sandstad.} \emph{The Black Hole Genesis}.
Physical Review D 97, 043504 (2018).
\code{doi.org/10.1088/1475-7516/2018/04/009}. The standard review of the
primordial black hole windows; Chapter~\ref{ch:pbh}.
\item \textbf{Carr and Kuhnel.} \emph{Primordial black hole dark matter and
the LIGO/Virgo merger events}. arXiv:2009.02424.
\item \textbf{Primordial black holes as dark matter}. arXiv:2312.17217.
\item \textbf{Carr, Clesse and Fattorini.} \emph{Ghostly galaxies}. Physical
Review D 101, 083012 (2020). \code{doi.org/10.1103/PhysRevD.101.083012}.
\item \textbf{Primordial black holes and the next generation of
gravitational-wave detectors}. arXiv:2602.15026.
\item \textbf{Cusp-core profile formation}. arXiv:1502.03456.
\item \textbf{Superheavy dark matter and beam dumps}. arXiv:2210.09313.
\item \textbf{Arbey and Mahmoudi.} \emph{Dark matter and the early
universe}. Progress in Particle and Nuclear Physics 95 (2017).
\code{doi.org/10.1142/S0217751X19300138}.
\end{enumerate}

\vspace{8pt}
{\sffamily\bfseries\small The gravitational evidence, and the alternatives}

\begin{enumerate}[leftmargin=1.4em,itemsep=4pt,start=52]
\item \textbf{Babcock.} \emph{The Rotation of the Galaxy NGC 3190}.
Astrophysical Journal 99, 41 (1939).
\code{doi.org/10.5479/ads/bib/1939licob.19.41b}. The first measurement of
flat rotation curves.
\item \textbf{Rubin and Ford.} \emph{Dark Matter}. Astrophysical Journal
159, L35 (1970). \code{doi.org/10.1086/150317}.
\item \textbf{Clowe et al.} \emph{A Direct Empirical Proof of the Existence of
Dark Matter}. Astrophysical Journal 648, L21 (2006).
\code{doi.org/10.1086/508162}. The Bullet Cluster.
\item \emph{How Dark Matter Came to Matter}. Nature Astronomy 1, 0009 (2017).
\code{doi.org/10.1038/s41550-017-0059}.
\item \textbf{van den Bergh.} \emph{The Early History of Dark Matter}.
Astronomical Journal 136, 1129 (2008).
\code{doi.org/10.1086/316369}.
\item \textbf{The Bullet Cluster and the limits of baryonic explanations.}
Chapter~\ref{ch:bullet}; the baryon budget is discussed in
\code{doi.org/10.1086/508162} and the cluster baryon fraction in the
review of Ref.~35.
\end{enumerate}

\vspace{8pt}
{\sffamily\bfseries\small The Solar System, the wind, and harvesting}

\begin{enumerate}[leftmargin=1.4em,itemsep=4pt,start=58]
\item \textbf{Dark matter in the Solar System}. arXiv:0806.3767.
\item \textbf{Capture of dark matter in the Solar System}. arXiv:0906.2480.
\item \textbf{Dark matter in the Solar System: the halo picture}.
arXiv:2504.07232.
\item \textbf{Dark matter wind}. Journal of High Energy Physics 07 (2024)
136. \code{doi.org/10.1007/JHEP07(2024)136}.
\item \textbf{Sikivie.} \emph{Dark Matter as a Dark Current}. arXiv:2107.14300.
\item \textbf{Dark matter power in quantum devices}. Physical Review Letters
132, 121801 (2024). \code{doi.org/10.1103/PhysRevLett.132.121801}.
\item \textbf{LOFAR and the solar corona}. Nature Communications (2024).
\code{doi.org/10.1038/s41467-024-45033-4}.
\end{enumerate}

\vspace{8pt}
{\sffamily\bfseries\small Strategy}

\begin{enumerate}[leftmargin=1.4em,itemsep=4pt,start=65]
\item \textbf{Particle Physics Opportunity Assessment.} arXiv:2104.07634.
The 2023 assessment; the portfolio argument; Chapter~\ref{ch:portfolio}.
\item \textbf{Snowmass 2021.} \emph{Cosmology Intertwined}. arXiv:2211.07027.
\item \textbf{Snowmass 2021.} \emph{Dark Matter Complementarity}. arXiv:2210.01770.
\item \textbf{The CF1 decision tree}. arXiv:1310.8327.
\end{enumerate}

\vspace{8pt}
\noindent{\small\color{txtdim}\sffamily
Excluded sources. Three claims that circulated widely were excluded from this
book because no primary record could be located for them: a reported 2025
NASA announcement of the first detection of dark matter in the Solar System;
a non-peer-reviewed candidate ranking; and a monograph titled \emph{Dark
Matter and Power} that could not be resolved in any bibliographic
registry. The source ledger accompanying this book records the attempts and
the reasons.\par}

% ---- Further Reading -------------------------------------------------------
\chapter*{Further Reading}
\addcontentsline{toc}{chapter}{Further Reading}
\markboth{Further Reading}{Further Reading}
\thispagestyle{plain}

\noindent\small
For readers who want to go further than this book goes, in roughly the order
in which the subject is best learned.

\begin{itemize}
\item \textbf{The cosmological measurement.} Planck Collaboration,
\emph{Planck 2018 results. VI. Cosmological parameters} (2020), and the
DESI DR2 papers on baryon acoustic oscillations (2025). These two together
are the empirical foundation of everything in Part~I.
\item \textbf{The experimental frontier.} The LZ, XENONnT and PandaX
publications listed in the Notes and References, read in order of
publication, and the Particle Data Group review of dark matter searches for
the summary and for the current state of the exclusions.
\item \textbf{The candidates.} The review of dark matter (arXiv:1908.00612)
is the best single survey of the candidate space; the axion case is set out
in the ADMX papers and in the Particle Data Group review; the primordial
black hole case is set out by Carr and Kuhnel.
\item \textbf{The small-scale problems.} The cusp-core, missing satellites
and too-big-to-fail literature is best approached through the reviews of
Bullock and Boylan-Kolchin, and the soliton question through De Martino and
through the wave-like lensing work.
\item \textbf{Strategy.} The Particle Physics Opportunity Assessment and the
two Snowmass white papers are the best statements of where the field should
go, and both are written to be read by non-specialists.
\item \textbf{The history, for context only.} Van den Bergh's history of the
field, and the Nature Astronomy article on how the concept of dark matter
came to matter. These are useful for the framing and are not needed for any
of the arguments in this book.
\end{itemize}

% ---- Index ------------------------------------------------------------------
\chapter*{Index}
\addcontentsline{toc}{chapter}{Index}
\markboth{Index}{Index}
\thispagestyle{plain}
\printindex

% ---- Acknowledgments -------------------------------------------------------
\chapter*{Acknowledgments}
\addcontentsline{toc}{chapter}{Acknowledgments}
\markboth{Acknowledgments}{Acknowledgments}
\thispagestyle{plain}

This book exists because a great deal of careful work was published in the
open, and because the question it addresses is more interesting than any of
the partial answers available.

The people who built the cosmological measurement over more than half a
century are owed more than a footnote each. Their results are the terrain on
which everything in Part~I rests. The instrument collaborations --- Planck,
DESI, DES, the Rubin Observatory and the rest --- are the reason the
measurements in this book have the precision they have, and the people who
run the detectors through the night, keep the xenon cold, and publish their
data in public whether the result is positive or negative, are the ones who
made the field's culture what it is.

The collaborations whose null results are the subject of Chapter~\ref{ch:direct}
deserve particular thanks, because publishing a limit is not a result that
advances a career in the way a discovery does, and the entire exclusion
frontier exists because of that work.

Several claims discussed in this book were made by groups whose claims did
not survive scrutiny. The record in Chapter~\ref{ch:indirect} treats those
cases as errors in statistical and instrumental reasoning rather than as
misconduct, and the thanks are owed to the groups that performed the
non-confirmations and the re-analyses, which cost them credit and bought the
field its standards.

The verification work behind this book --- resolving every identifier
against a bibliographic registry, checking each resolved title against the
paper it was intended to cite, and removing the claims that could not be
traced --- was done with the assistance of the tooling in the accompanying
source materials. The corrections log records the places where that process
changed what the book says, including several that would otherwise have made
it wrong.

\begin{flushright}
\emph{Gerald Enrique Nelson Mc Kenzie}\\
Tico Internet LLC\\
22 November 2026
\end{flushright}

% ---- About the author ------------------------------------------------------
\chapter*{About the Author}
\addcontentsline{toc}{chapter}{About the Author}
\markboth{About the Author}{About the Author}
\thispagestyle{plain}

\noindent\textbf{Gerald Enrique Nelson Mc Kenzie}\par
\vspace{6pt}

\noindent\small
Gerald Enrique Nelson Mc Kenzie is the author of \emph{Uniworld: A
Relational Monism of Inseparable Frames} (Tico Internet LLC, 2026 First
Edition, Apache License 2.0, DOI 10.5281/zenodo.24119718) and of
\emph{Dark Matter: Its True Origins, Whether We Have Already Tapped It, and
Why the Answer Matters} (Tico Internet LLC, 2026 First Edition, Apache
License 2.0, DOI 10.5281/zenodo.24435288).

He writes about research fields at the moment they are being decided rather
than after they have settled, and about the evidence rather than the
headlines. His interest in dark matter is practical. A field whose central
quantity is known to two per cent and whose central object has never been
observed is a field in which the quality of the argument is the only thing
that distinguishes progress from noise, and the argument has to be made in
public, against sources that can be checked.

He publishes under the imprint Tico Internet LLC.

\vfill
\noindent{\sffamily\small\color{txtdim}\rule{\textwidth}{0.4pt}}\par
\noindent{\sffamily\footnotesize\color{txtdim}
\emph{Corrections.} Corrections to this edition are welcomed and will be
carried in a later printing. The corrections log accompanying the research
materials for this title records, in the open, the places where the writer's
initial assumptions were revised during preparation.\par}

% ---- Colophon -------------------------------------------------------------
\clearpage
\thispagestyle{empty}
\vspace*{1.0in}
\begin{center}
\begin{minipage}{0.92\textwidth}
{\sffamily\fontsize{15}{18}\selectfont\bfseries\color{txt} \bookTitle\par}
\vspace{4pt}
{\sffamily\small\color{txtdim} \bookSubtitle\par}
\vspace{18pt}
{\color{accent}\rule{\textwidth}{1.6pt}}\par
\vspace{14pt}
{\small\color{txt}
\noindent\textbf{Author}\quad \bookAuthor\par
\vspace{4pt}
\noindent\textbf{Publisher}\quad \bookImprint\par
\vspace{4pt}
\noindent\textbf{Edition}\quad \bookEdition\par
\vspace{4pt}
\noindent\textbf{DOI}\quad \href{https://doi.org/\bookDOI}{\texttt{\bookDOI}}\par
\vspace{4pt}
\noindent\textbf{Licence}\quad \bookLicense\par
\vspace{4pt}
\noindent\textbf{Source materials}\quad
\texttt{github.com/lordxmen2k/Dark-Matter-True-Origins}. The book
sources, the build scripts, the figures and the research ledger behind this
book are published there.\par
\vspace{14pt}
\noindent\textbf{Subject headings}\quad
Cosmology --- Dark matter; Astrophysics --- Cosmology; Particle physics ---
Dark matter; Direct detection (Astronomy); Axions (Physics); Black holes
(Astronomy); Cosmological parameters.\par
\vspace{14pt}
\noindent\textbf{Errata}\quad
No errata have been recorded for this edition.\par
\vspace{14pt}
\noindent\textbf{Extent}\quad
\bookPages{} numbered pages, with full-bleed front and back
covers;\ \bookWords{} words.\par
\vspace{14pt}
\noindent\rule{\textwidth}{0.4pt}\par
\vspace{10pt}
\noindent\textbf{Signed}\quad
\emph{Gerald Enrique Nelson Mc Kenzie}\par
\vspace{4pt}
\noindent Tico Internet LLC\par
\vspace{4pt}
\noindent \bookFactsDate\par
}
\end{minipage}
\end{center}
